\documentclass{article}
\usepackage{fleqn}
\usepackage{epsf}
\usepackage{aima2e-slides}

\begin{document}

\begin{huge}
\titleslide{Statistical learning}{Chapter 20, Sections 1--3}

\sf

%%%%%%%%%%%% Slide %%%%%%%%%%%%%%%%%%%%%%%%%%%%%%%%%%%%%%%%%%%%%%%%%%%
\heading{Phác thảo}

\blob Học Bayes

\blob Tối đa {\em a posteriori} và khả năng học tập tối đa

\blob Học mạng Bayes\nl
-- Học tham số ML với dữ liệu hoàn chỉnh\nl
-- hồi quy tuyến tính


%%%%%%%%%%%% Slide %%%%%%%%%%%%%%%%%%%%%%%%%%%%%%%%%%%%%%%%%%%%%%%%%%%
\heading{Học Bayesian đầy đủ}

Xem việc học dưới dạng cập nhật Bayesian của phân bố xác suất\\
trên không gian giả thuyết \defn{}

\mat{$H$} là biến giả thuyết, giá trị \mat{$h_1,h_2,\ldots$}, trước \mat{$\pv(H)$}

\mat{$j$}quan sát thứ \mat{$d_j$} đưa ra kết quả của biến ngẫu nhiên \mat{$D_j$}\\
dữ liệu huấn luyện \mat{$\d \eq d_1,\ldots,d_{\Ncount}$}


Với dữ liệu cho đến nay, mỗi giả thuyết có xác suất sau:
\mat{\[
  P(h_i|\d) = \alpha P(\d|h_i) P(h_i)
\]}%
trong đó \mat{$P(\d|h_i)$} được gọi là \defn{khả năng}

Các dự đoán sử dụng mức trung bình có trọng số khả năng đối với các giả thuyết:
\mat{\[
  \pv(X|\d) = \mysum_i\ \pv(X|\d,h_i) P(h_i|\d) = \mysum_i\ \pv(X|h_i) P(h_i|\d)
\]}%
Không cần phải chọn một giả thuyết dự đoán tốt nhất!


%%%%%%%%%%%% Slide %%%%%%%%%%%%%%%%%%%%%%%%%%%%%%%%%%%%%%%%%%%%%%%%%%%
\heading{Ví dụ}

Giả sử có năm loại túi kẹo:\al
10\% là \mat{$h_1$}: 100\% kẹo anh đào\al
20\% là \mat{$h_2$}: 75\% kẹo anh đào + 25\% kẹo chanh\al
40\% là \mat{$h_3$}: 50\% kẹo anh đào + 50\% kẹo chanh\al
20\% là \mat{$h_4$}: 25\% kẹo anh đào + 75\% kẹo chanh\al
10\% là \mat{$h_5$}: 100\% kẹo chanh

\vspace*{0.2in}

\epsfxsize=0,6\textwidth
\fig{\file{figures}{candy-kinds.ps}}

Sau đó, chúng tôi quan sát những viên kẹo được rút ra từ một số túi: \epsfxsize=0.3\textwidth \epsffile{\file{figures}{candy-obs.ps}}

Đó là loại túi gì? Kẹo tiếp theo sẽ có hương vị gì?


%%%%%%%%%%%% Slide %%%%%%%%%%%%%%%%%%%%%%%%%%%%%%%%%%%%%%%%%%%%%%%%%%%
\heading{Xác suất sau của giả thuyết}

\vspace*{0.2in}

\epsfxsize=0.9\textwidth
\fig{\file{graphs}{limes-bayes-post.eps}}




%%%%%%%%%%%% Slide %%%%%%%%%%%%%%%%%%%%%%%%%%%%%%%%%%%%%%%%%%%%%%%%%%%
\heading{Xác suất dự đoán}

\vspace*{0.2in}

\epsfxsize=0.9\textwidth
\fig{\file{graphs}{limes-bayes-pred.eps}}


%%%%%%%%%%%% Slide %%%%%%%%%%%%%%%%%%%%%%%%%%%%%%%%%%%%%%%%%%%%%%%%%%%
\heading{Xấp xỉ MAP}

Tổng hợp không gian giả thuyết thường khó thực hiện \\
(ví dụ: 18,446,744,073,709,551,616 hàm Boolean của 6 thuộc tính)

\defn{ Học tối đa sau } (MAP): chọn \mat{$h_{{\rm MAP}}$} tối đa hóa \mat{$P(h_i|\d)$}

Tức là tối đa hóa \mat{$P(\d|h_i)P(h_i)$} hoặc \mat{$\log P(\d|h_i) + \log P(h_i)$}

Thuật ngữ nhật ký có thể được xem là (phủ định của)\al
\note{bit để mã hóa dữ liệu cho giả thuyết} + \note{bit để mã hóa giả thuyết} \\
Đây là ý tưởng cơ bản của việc học \defn{độ dài mô tả tối thiểu} (MDL)

Đối với các giả thuyết xác định, \mat{$P(\d|h_i)$} là 1 nếu nhất quán, 0 nếu không \nl
$\implies$ MAP = giả thuyết nhất quán đơn giản nhất (cf.~science)



%%%%%%%%%%%% Slide %%%%%%%%%%%%%%%%%%%%%%%%%%%%%%%%%%%%%%%%%%%%%%%%%%%
\heading{Xấp xỉ ML}

Đối với các tập dữ liệu lớn, trước trở nên không liên quan

\defn{Khả năng tối đa} (ML): chọn \mat{$h_{{\rm ML}}$} tối đa hóa \mat{$P(\d|h_i)$}

Tức là, chỉ cần lấy dữ liệu phù hợp nhất; giống hệt với MAP cho đồng phục trước\\
(điều này hợp lý nếu tất cả các giả thuyết đều có cùng độ phức tạp)

ML là phương pháp học thống kê ``chuẩn'' (không phải Bayesian)


%%%%%%%%%%%% Slide %%%%%%%%%%%%%%%%%%%%%%%%%%%%%%%%%%%%%%%%%%%%%%%%%%%
\heading{Học tham số ML trong mạng Bayes}

Túi từ nhà sản xuất mới; phần \mat{$\theta$} của kẹo anh đào?\hfill\epsfxsize=1.25in\raisebox{-1.25in[0pt][0pt]{\epsffile{\file{figures}{ml-network1.ps}}} \\
Bất kỳ \mat{$\theta$} nào cũng có thể xảy ra: sự liên tục của các giả thuyết \mat{$h_{\theta}$} \\
\mat{$\theta$} là tham số \defn{} cho dòng mô hình đơn giản (\note{binomial}) này 

Giả sử chúng ta mở gói \(\Ncount\) kẹo, \(c\) quả anh đào và \(\ell\eq \Ncount-c\) chanh\\
Đây là các quan sát \defn{i.i.d.} (độc lập, phân bố giống hệt nhau), vì vậy
\mat{\[
  P(\data |h_{\theta}) = \prod_{j\eq 1}^{\Ncount} P(\datum_j|h_{\theta}) =
  \theta^c\cdot (1-\theta)^{\ell}
\]}%
Tối đa hóa w.r.t.~\mat{$\theta$}---điều này dễ dàng hơn đối với \defn{log-likelihood}:
\mat{\begin{eqnarray*}
  L(\data |h_{\theta}) &=& \log P(\data |h_{\theta}) = \sum_{j\eq 1}^{\Ncount} \log P(\datum_j|h_{\theta}) =
  c\log\theta +  \ell\log(1-\theta) \\
  \frac{dL(\data |h_{\theta})}{d\theta} &=& \frac{c}{\theta} -
  \frac{\ell}{1-\theta} = 0 \qquad \implies \quad \theta =
  \frac{c}{c+\ell} = \frac{c}{\Ncount}
\end{eqnarray*}}%
Có vẻ hợp lý nhưng gây ra vấn đề với số 0!








%%%%%%%%%%%% Slide %%%%%%%%%%%%%%%%%%%%%%%%%%%%%%%%%%%%%%%%%%%%%%%%%%%
\heading{Nhiều tham số}

Giấy gói màu đỏ/xanh phụ thuộc xác suất vào hương vị: \hspace*{0.75in}\epsfxsize=2.25in\raisebox{-2.5in[0pt][0pt]{\epsffile{\file{figures}{ml-network2.ps}}} 

Ví dụ: khả năng xảy ra với kẹo anh đào trong giấy gói màu xanh lá cây:
\mat{\begin{eqnarray*}
\lefteqn{P(\J{F}\eq \J{cherry},\J{W}\eq \J{green}|h_{\theta,\theta_1,\theta_2})}\\
 & = & P(\J{F}\eq \J{cherry}|h_{\theta,\theta_1,\theta_2})P(\J{W}\eq
 \J{green}|\J{F}\eq \J{cherry},h_{\theta,\theta_1,\theta_2}) \\
 & = & \theta \cdot (1-\theta_1)
\end{eqnarray*}}%
\(\Ncount\) kẹo, \(r_c\) kẹo anh đào bọc đỏ, v.v.:
\mat{\begin{eqnarray*}
 P(\data|h_{\theta,\theta_1,\theta_2}) &=&
  \theta^c (1-\theta)^{\ell} \cdot \theta_1^{r_c}(1-\theta_1)^{g_c}
  \cdot \theta_2^{r_{\ell}}(1-\theta_2)^{g_{\ell}}
\end{eqnarray*}}%
\mat{\begin{eqnarray*}
L &=&  [c\log \theta + \ell\log (1-\theta) ]  \\
  &+&   [r_c\log \theta_1 + g_c\log(1-\theta_1) ]  \\
  &+&   [r_{\ell}\log \theta_2 + g_{\ell}\log(1-\theta_2) ]
\end{eqnarray*}}%

%%%%%%%%%%%% Slide %%%%%%%%%%%%%%%%%%%%%%%%%%%%%%%%%%%%%%%%%%%%%%%%%%%
\heading{Tiếp theo nhiều tham số.}

Đạo hàm của \mat{$L$} chỉ chứa tham số liên quan:
\mat{\[\begin{array}{rclcl}
\displaystyle \frac{\partial L}{\partial\theta} &=& \displaystyle \frac{c}{\theta} - \displaystyle\frac{\ell}{1-\theta} = 0                  & \qquad \implies & \theta = \displaystyle \frac{c}{c+\ell} 
\end{array}\]}%
\mat{\[\begin{array}{rclcl}
\displaystyle \frac{\partial L}{\partial\theta_1} &=& \displaystyle\frac{r_c}{\theta_1} - \displaystyle\frac{g_c}{1-\theta_1} = 0           & \qquad \implies & \theta_1 = \displaystyle\frac{r_c}{r_c+g_c} 
\end{array}\]}%
\mat{\[\begin{array}{rclcl}
\displaystyle \frac{\partial L}{\partial\theta_2} &=& \displaystyle\frac{r_{\ell}}{\theta_2} - \displaystyle\frac{g_{\ell}}{1-\theta_2} = 0 & \qquad \implies & \theta_2 = \displaystyle\frac{r_{\ell}}{r_{\ell}+g_{\ell}}
\end{array}\]}%
Với \defn{dữ liệu đầy đủ}, các tham số \emph{ có thể được học riêng }




%%%%%%%%%%%% Slide %%%%%%%%%%%%%%%%%%%%%%%%%%%%%%%%%%%%%%%%%%%%%%%%%%%
\heading{Ví dụ: mô hình Gaussian tuyến tính}

\twofig{\file{graphs}{regression-model.eps}}{\file{graphs}{regression-sample.eps}}

Tối đa hóa \mat{$\displaystyle P(y|x) = \frac{1}{\sqrt{2 \pi} \sigma} e^{-\frac{(y-(\theta_1 x+\theta_2))^2}{2\sigma^2}}$} w.r.t. \mat{$\theta_1$}, \mat{$\theta_2$}

= giảm thiểu \mat{$\displaystyle E = \sum_{j\eq 1}^{\Ncount} (y_j-(\theta_1 x_j+\theta_2))^2$}

Nghĩa là, giảm thiểu tổng sai số bình phương sẽ mang lại giải pháp ML\\
để phù hợp tuyến tính \emph{giả sử nhiễu Gaussian có phương sai cố định }



%%%%%%%%%%%% Slide %%%%%%%%%%%%%%%%%%%%%%%%%%%%%%%%%%%%%%%%%%%%%%%%%%%
\heading{Tóm tắt}

Học tập Bayesian đầy đủ đưa ra dự đoán tốt nhất có thể nhưng khó thực hiện

Học MAP cân bằng độ phức tạp với độ chính xác trên dữ liệu đào tạo 

Khả năng tối đa giả định thống nhất trước đó, OK đối với các tập dữ liệu lớn

1. Chọn một nhóm mô hình được tham số hóa để mô tả dữ liệu\\
\note{{\em yêu cầu cái nhìn sâu sắc đáng kể và đôi khi là các mô hình mới}

2. Viết ra khả năng xảy ra của dữ liệu dưới dạng hàm của các tham số\\
\note{{\em có thể yêu cầu tính tổng các biến ẩn, tức là suy luận}

3. Viết đạo hàm của log khả năng w.r.t. mỗi tham số

4. Tìm các giá trị tham số sao cho đạo hàm bằng 0\\
\note{{\em có thể khó/không thể; trợ giúp về kỹ thuật tối ưu hóa hiện đại}



\end{huge}
\end{document}


